\section{Simetria e Assimetria}

Entre as motivações para se estudar a assimetria ou simetria dos dados estão a constação
de
\begin{itemize}
\item Concentração dos dados em uma dada região de valores.
\item Presença de caldas particularmente longas.
\item Afastamento dos dados do centro.
\item Validade de hipóteses sobre o modelo adotado para a distribuição.
\end{itemize}

Quando os dados possuem distribuição simétrica, algumas vantagens analíticas e estatísticas
surgem.
Por exemplo, a média se torna representativa para medir centralidade, pois essa localiza-se
justamente no centro da distribuição, de sorte que a variância -- e o desvio-padrão --
também passa a ser representativa para medir variabilidade, já que os dados estão
centrados na média.
Quando os dados são assimétricos, medidas de resumo como mediana e qartis são preferíveis
por serem menos sensíveis à forma da distribuição.
Algumas técnicas de inferência estatística, principalmente aquelas que se baseiam no uso
da distribuição normal, a serem apresentadas adiante, também se aproveitam da estrutura
simétrica da distribuição, facilitando a produção de algumas análises.

Seja $X$ um atributo.
Para medir sua assimetria, são geralmente usados
\begin{itemize}
\item Primeiro Coeficiente de Assimetria de Pearson -- usado quando a distribuição é
unimodal, e a moda pode ser bem definida.
\[
   \frac{\overline x - \mo X}{\sigma^3}
\]

\item Segundo Coeficiente de Assimetria de Pearson -- usado quando a distribuição não
é unimodal ou a moda não pode ser bem definida, dada por
\[
     \frac{3\,(\overline x - \med X)}{\sigma^3}
\]

\item Coeficiente Quatílico de Assimetria -- utiliza os quartis da distribuição, partindo
do princípio que o primeiro e terceiro quartil são equidistantes da mediana
\[
    \frac{q_1 + q_2 - 2\, \med X}{q_3 - q_1}
\]

\item Coeficiente de Assimetria de Fisher -- um estimador baseado no terceiro momento
central multiplicado por um fator que remove o viés do tamanho da amostra
\[
    \frac{\sqrt{n(n-1)}}{n-2} \cdot \frac{1}{n\sigma^3} \cdot \sum_{i=1}^n (x_i - \overline x)^3
\]
\end{itemize}
A depender da medida $c$ utilizada, ela pode ter seus valores interpretados como
\begin{itemize}
\item $c > 0$ -- assimetria positiva, indicando que $\overline X > \mo X$, também
chamada de assimetria à esquerda, pois os dados possuem uma calda mais longa nessa
direção enquanto que seus valores concentram-se mais à direita.

\item $c < 0$ -- assimetria negativa, e provavelmente temso $\overline X < \mo X$, o que
é chamado de assimetria à direita, direção na qual se estende a calda da distribuição,
e os dados concentram-se em valores mais baixos.

\item $c = 0$ -- ausência de assimetria, levando a hipótese de que $\overline X = \mo X$
e que a distribuição é possivelmente simétrica.
\end{itemize}

Uma ferramente visual para medir assimetria de dados são os \textbf{gráficos de simetria}.
As estatísticas de ordem são usadas para computar pares $(u_i, v_i)$ tais que
\[
    u_i = \med X - x_i \quad \text{e} \quad v_i = x_{n-i+1} - \med X.
\]
Junto desses pares, é mostrado também a reta $v = u$.
Se os pontos se alinharem a reta, temos uma evidência visual de que os dados tem pouca
assimetria ou são simétricos, enquanto que, caso alguns pontos se desviem das extreidades
dela, então teremos uma visualização que indica assimetria.
